\documentclass[11pt]{article}
\usepackage[margin=0.85in]{geometry}
\usepackage{amsmath}
\usepackage{enumitem}
\usepackage{xcolor}
\usepackage{parskip}
\usepackage{titlesec}
\definecolor{ink}{HTML}{1F2733}
\definecolor{accent}{HTML}{C05621}
\definecolor{teal}{HTML}{16545C}
\titleformat{\section}{\large\bfseries\color{teal}}{}{0em}{}
\titlespacing{\section}{0pt}{8pt}{2pt}
\setlist[enumerate]{leftmargin=1.4em,itemsep=3pt,topsep=2pt}
\color{ink}
\pagestyle{empty}

\begin{document}

{\LARGE\bfseries Pre-Hackathon Consultation --- Agenda}\\[2pt]
{\color{accent}\rule{\textwidth}{1.2pt}}\\[4pt]
\textbf{Project:} Spatial Degradation of Hepatocyte Zonation \quad\textbf{Team:} Roee \& Shira
\hfill \textbf{Course:} Comp.\ Genomics 76553

\smallskip
\noindent\textit{One line:} use Paper~2's healthy human-liver zonation atlas as a quantitative
``ruler'' to give Paper~1's spatially-blind disease hepatocytes a pseudo-zonation coordinate, then
measure how zonation collapses across MASLD stages (H1), which gene programs collapse (H2), and
whether de-zonated cells turn plastic (H3).

\smallskip
\noindent\fbox{\parbox{0.97\textwidth}{\footnotesize
\textbf{Background for the questions.} The ruler is a per-cell coordinate
$=\operatorname{mean}_z(\text{pericentral genes})-\operatorname{mean}_z(\text{periportal genes})$,
with the gene sets taken from \textbf{Paper~2's healthy zonation tables} (genes scored as
significantly zonated, assigned to a zone by their center-of-mass). We have four nested sets, all
from Paper~2: \textbf{full} ($\sim$1{,}600 transcriptome-wide zonated genes), \textbf{expanded}
($\sim$100 top-ranked), \textbf{core} ($\sim$13/8 hand-curated), and the \textbf{landmark} set ---
Paper 2's exact 20+20 hepatocyte landmark genes (from its published \texttt{-LM} lists), which Paper 2
itself uses to assign snRNA zonation. The \textbf{``validation battery''} $=$ the checks that the coordinate truly measures
position: recovery of canonical markers on healthy cells (positive control), within-cell
pericentral--periportal anti-correlation, internal coherence of each module, and split-half
reproducibility.}}

\section{1. Methodology \& biology (calls that could change direction)}
\begin{enumerate}
  \item \textbf{Signature strategy.} We lead with the \emph{full} transcriptome-wide program (the
  novelty: read position from the whole program, not a few markers), and report the smaller
  canonical sets as an audit/robustness check; the validation battery decides which leads. Is this
  appropriately novel without over-reaching?
  \item \textbf{Classifier ground truth.} We will train the zone classifier on Paper~2's
  \emph{true Visium-mapped} zonation labels (not a coordinate-derived label), so its per-cell entropy
  is an \emph{independent} de-zonation signal. Is that the right ground truth, and any pitfalls in the
  Visium\,$\rightarrow$\,snRNA label transfer?
  \item \textbf{Biological prior on the collapse.} Is there a known \emph{directional} pattern (e.g.\
  pericentral injury earlier in MASH) we should test as a specific hypothesis, rather than only a
  generic monotonic decline?
  \item \textbf{Is our metric ``de-zonation''?} We quantify it per donor as: coordinate spread
  narrowing, pericentral--periportal anti-correlation weakening toward~0, and classifier entropy
  rising. Do these capture ``loss of zonation'' as you would define it?
  \item \textbf{Plasticity axis (H3).} Are cholangiocyte / ductular-reaction markers
  (\textit{KRT7, KRT19, SOX9}) the right readout for hepatocyte plasticity in MASLD, or is there a
  more established de-differentiation signature?
\end{enumerate}

\section{2. Scope \& expectations}
\begin{enumerate}
  \item \textbf{What does top-mark work look like?} We aim to deliver \emph{all} of H1--H3 (plus a
  bonus enrichment) in two days. Given that, what most distinguishes an excellent project here ---
  depth, breadth, rigor of controls, or the write-up / talk?
  \item \textbf{Small $n$ at the extremes.} Cirrhosis ($n{=}4$) and end-stage ($n{=}5$) rest on few
  donors. Acceptable as a stated limitation with donor-level confidence intervals, or would you
  handle it differently?
\end{enumerate}

\vfill
{\color{teal}\rule{\textwidth}{0.4pt}}\\[2pt]
\footnotesize Unit of inference throughout: the \textbf{donor} ($\approx$47 patients), never the cell
($\approx$69k) --- to avoid pseudoreplication.

\end{document}
